% An elicited two-measure additive model; generic, not empirical preference data.
\begin{tikzpicture}[font=\small,>=Stealth,
 box/.style={draw,align=center,minimum height=10mm,text width=27mm},
 arr/.style={->,thick}]
 \node[box,fill=blue!8] (x1) at (0,1.1) {Measure score\\$x_1$};
 \node[box,fill=blue!8] (x2) at (0,-1.1) {Measure score\\$x_2$};
 \node[box] (f1) at (3.5,1.1) {Elicited value\\$f_1(x_1)\in[0,1]$};
 \node[box] (f2) at (3.5,-1.1) {Elicited value\\$f_2(x_2)\in[0,1]$};
 \node[box,fill=orange!12,text width=32mm] (v) at (7,0) {$V=w_1f_1+w_2f_2$\\$w_1,w_2\ge0$\\$w_1+w_2=1$};
 \draw[arr] (x1)--(f1); \draw[arr] (x2)--(f2);
 \draw[arr] (f1.east)--node[above,sloped] {$w_1$} (v.north west);
 \draw[arr] (f2.east)--node[below,sloped] {$w_2$} (v.south west);
 \node[align=center] at (3.5,-2.15) {Value increments and weights describe preferences;\\normalizing the numbers alone does not elicit them.};
\end{tikzpicture}
